% 基于教师提供的 Overleaf 模板：https://www.overleaf.com/read/hnsjbjyknbdj
\usepackage[left=3.17cm, right=3.17cm, top=2.74cm, bottom=2.74cm]{geometry}
\usepackage{amsmath}
\usepackage{graphicx,subfig}
\usepackage{float}
\usepackage{cite}
\usepackage{caption}
\usepackage{enumerate}
\usepackage{booktabs}
\newcommand{\tabincell}[2]{\begin{tabular}{@{}#1@{}}#2\end{tabular}}

% 显式配置中西文字体，避免依赖 ctex 按操作系统自动选择的字体集。
\usepackage{fontspec}
\usepackage{xeCJK}
\IfFileExists{/System/Library/Fonts/Supplemental/Songti.ttc}{%
  \setCJKmainfont{Songti.ttc}[Path=/System/Library/Fonts/Supplemental/,FontIndex=6]}{%
  \IfFontExistsTF{Noto Serif CJK SC}{\setCJKmainfont{Noto Serif CJK SC}}{\setCJKmainfont{SimSun}}}
\IfFileExists{/System/Library/Fonts/STHeiti Medium.ttc}{%
  \setCJKsansfont{STHeiti Medium.ttc}[Path=/System/Library/Fonts/,FontIndex=1]}{%
  \IfFontExistsTF{Noto Sans CJK SC}{\setCJKsansfont{Noto Sans CJK SC}}{\setCJKsansfont{Microsoft YaHei}}}
\setmonofont{Courier New}
\IfFileExists{/System/Library/Fonts/STHeiti Medium.ttc}{%
  \setCJKmonofont{STHeiti Medium.ttc}[Path=/System/Library/Fonts/,FontIndex=1]}{%
  \IfFontExistsTF{Noto Sans Mono CJK SC}{\setCJKmonofont{Noto Sans Mono CJK SC}}{\setCJKmonofont{Microsoft YaHei}}}
\IfFileExists{/System/Library/AssetsV2/com_apple_MobileAsset_Font8/88d6cc32a907955efa1d014207889413890573be.asset/AssetData/Kaiti.ttc}{%
  \newCJKfontfamily\kaishu{Kaiti.ttc}[Path=/System/Library/AssetsV2/com_apple_MobileAsset_Font8/88d6cc32a907955efa1d014207889413890573be.asset/AssetData/,FontIndex=0]}{%
  \IfFontExistsTF{Noto Serif CJK SC}{\newCJKfontfamily\kaishu{Noto Serif CJK SC}}{%
    \newCJKfontfamily\kaishu{SimSun}}}
\newcommand{\yihao}{\fontsize{26pt}{36pt}\selectfont}
\newcommand{\erhao}{\fontsize{22pt}{28pt}\selectfont}
\newcommand{\xiaoer}{\fontsize{18pt}{18pt}\selectfont}
\newcommand{\sanhao}{\fontsize{16pt}{24pt}\selectfont}
\newcommand{\xiaosan}{\fontsize{15pt}{22pt}\selectfont}
\newcommand{\sihao}{\fontsize{14pt}{21pt}\selectfont}
\newcommand{\banxiaosi}{\fontsize{13pt}{19.5pt}\selectfont}
\newcommand{\xiaosi}{\fontsize{12pt}{18pt}\selectfont}
\newcommand{\dawuhao}{\fontsize{11pt}{11pt}\selectfont}
\newcommand{\wuhao}{\fontsize{10.5pt}{15.75pt}\selectfont}

% 模板章节、页眉页脚与封面横线设置。
\usepackage{ctexcap}
\CTEXsetup[name={,、},number={ \chinese{section}}]{section}
\CTEXsetup[name={（,）},number={\chinese{subsection}}]{subsection}
\CTEXsetup[name={,.},number={\arabic{subsubsection}}]{subsubsection}
\usepackage{fancyhdr}
\pagestyle{fancy}
\lhead{\kaishu \leftmark}
\rhead{\kaishu 编译原理上机作业调研报告}
\lfoot{}
\cfoot{\thepage}
\rfoot{}
\renewcommand{\headrulewidth}{0.1pt}
\renewcommand{\footrulewidth}{0pt}
\newcommand{\HRule}{\rule{\linewidth}{0.5mm}}
\newcommand{\HRulegrossa}{\rule{\linewidth}{1.2mm}}

% 模板代码、高亮、超链接与水印设置。
\usepackage{xcolor}
\usepackage{listings}
\lstset{
  breaklines=true,
  basicstyle=\ttfamily\small,
  keywordstyle=\color{blue!70}\bfseries,
  commentstyle=\color{red!50!green!50!blue!50},
  stringstyle=\ttfamily,
  extendedchars=false,
  linewidth=\textwidth,
  numbers=left,
  numberstyle=\tiny\color{blue!50},
  frame=trbl,
  rulesepcolor=\color{red!20!green!20!blue!20},
  numbersep=8pt,
  showstringspaces=false,
  columns=fullflexible,
  keepspaces=true,
  tabsize=4,
  captionpos=b,
  backgroundcolor=\color{white},
  rulecolor=\color{black}
}
\lstdefinelanguage{bash}{
  morekeywords={cd,ls,pwd,grep,find,head,tail,echo,file,nm,objdump,readelf,which,type,export,rm,cp,make},
  sensitive=true,
  morecomment=[l]{\#},
  morestring=[b]"
}
\lstdefinelanguage{RiscVAsm}{
  morekeywords={li,mv,j,call,ret,bnez,blez,bge,slli,add,lw,sw,remw,addw,addi,addiw,sd,ld},
  sensitive=true,
  morecomment=[l]{\#}
}
\usepackage[colorlinks,linkcolor=black,anchorcolor=blue]{hyperref}
\usepackage{tikz}
\usepackage{eso-pic}
\newcommand{\watermark}[3]{\AddToShipoutPictureBG{\parbox[b][\paperheight]{\paperwidth}{\vfill\centering\tikz[remember picture, overlay]\node [rotate = #1, scale = #2] at (current page.center){\textcolor{gray!80!cyan!30!magenta!30}{#3}};\vfill}}}

% 与现有实验内容对应的表格和排版支持。
\usepackage{tabularx}
\usepackage{array}
\usepackage{longtable}
\usepackage{enumitem}
\usepackage{setspace}
\setstretch{1.35}
\setlength{\parindent}{2em}
\setlength{\parskip}{0pt}
\captionsetup{font=small,labelsep=space}
\renewcommand{\figurename}{图}
\renewcommand{\tablename}{表}
\renewcommand{\lstlistingname}{代码}
\renewcommand{\refname}{参考文献}
\newcommand{\labfigure}[3]{%
\begin{figure}[H]
  \centering
  \includegraphics[width=0.92\linewidth,height=0.72\textheight,keepaspectratio]{#1}
  \caption{#2}
  \label{#3}
\end{figure}}
\newcommand{\labfigurecompact}[4]{%
\begin{figure}[H]
  \centering
  \includegraphics[width=#4\linewidth,height=0.72\textheight,keepaspectratio]{#1}
  \caption{#2}
  \label{#3}
\end{figure}}\setlength{\headheight}{14pt}
\emergencystretch=3em
